\section{Grundlagen}
\label{sec:fundamentals}

Dieser Abschnitt führt alle Konzepte ein, die für das Verständnis der Methodik (\S\ref{sec:methodology}) und der Ergebnisse (\S\ref{sec:results}) notwendig sind. Er richtet sich an Leser mit ML-Hintergrund, die nicht notwendigerweise mit Quantisierung oder Krümmungsanalyse vertraut sind.

\subsection{Quantisierung neuronaler Netze}
\label{sec:fund-quant}

\textbf{Grundidee.} Ein trainiertes neuronales Netz speichert seine Gewichte und berechnet seine Aktivierungen typischerweise als 32-Bit-Gleitkommazahlen (FP32). Quantisierung bildet diese kontinuierlichen Werte auf eine kleinere, diskrete Menge von Repräsentanten ab, üblicherweise mit deutlich weniger Bits pro Wert. Das reduziert den Speicherbedarf direkt proportional zur Bit-Breite und kann, wenn die Zielhardware quantisierte Arithmetik-Einheiten besitzt, auch die Rechenzeit senken.
\par\medskip

\textbf{PTQ vs.\ QAT.} Es gibt zwei grundsätzliche Strategien, ein Netz zu quantisieren. Bei der Post-Training-Quantisierung (PTQ) wird ein bereits fertig trainiertes FP32-Modell nachträglich quantisiert, ohne weiteres Training. PTQ ist schnell und benötigt keine Trainingsdaten in großem Umfang, produziert aber typischerweise größere Genauigkeitsverluste, da das Modell nie gelernt hat, mit dem Quantisierungsfehler umzugehen. Beim Quantization-Aware-Training (QAT) wird der Quantisierungseffekt bereits während des Trainings simuliert, sodass das Modell seine Gewichte an die Quantisierung anpassen kann. QAT erzielt in der Regel höhere Genauigkeit, benötigt aber zusätzliches Training.
\par\medskip

\textbf{Fake-Quantization und Straight-Through-Estimator.} Sowohl bei der PTQ-Kalibrierung als auch bei QAT wird der Quantisierungseffekt üblicherweise durch eine Fake-Quantization simuliert. Dabei wird ein Wert quantisiert und sofort wieder in Gleitkomma zurückkonvertiert (oft als \enquote{QDQ} bezeichnet), sodass der Rest der Berechnung weiterhin in FP32 abläuft, den Quantisierungsfehler aber dennoch \enquote{sieht}. Das eigentliche Problem liegt in der Rundungsfunktion in der Mitte dieser Operation, deren Ableitung fast überall null ist und die den Gradientenfluss beim Backpropagation-Training blockieren würde. Der Straight-Through-Estimator (STE, \cite{bengio_ste_2013}) löst dies, indem er im Vorwärtsdurchlauf die tatsächliche, nicht-differenzierbare Quantisierung anwendet, im Rückwärtsdurchlauf jedoch so verfährt, als wäre die Quantisierungsfunktion die Identität. Der Gradient wird also unverändert durchgereicht. Diese Relaxierung ist für dieses Projekt relevant, da sie auch bestimmt, welche Krümmung auf quantisierten Modellen tatsächlich gemessen wird (\S\ref{sec:fund-hessian}).
\par\medskip

\textbf{Power-of-Two-Quantisierung.} Während die uniforme Quantisierung die Werte auf ein gleichmäßig verteiltes Gitter abbildet (etwa $256$ äquidistante Stufen bei 8 Bit), rundet die PoT-Quantisierung jeden Gewichtswert auf die nächste Zweierpotenz innerhalb eines Exponentenbereichs \citep{miyashita_pot_2016}. Der entscheidende Vorteil ist arithmetischer Natur, da eine Multiplikation $x \cdot 2^k$ auf Digitalhardware einem Bitshift um $k$ Stellen entspricht. Dieser ist deutlich billiger als eine allgemeine Multiplikation und ohne dedizierten Multiplizierer realisierbar \citep{przewlocka_rus_pot_2022}. Der Nachteil ist ein nicht-uniformes Fehlerverhalten. Da das PoT-Gitter logarithmisch ist, ist die Gitterauflösung nahe null sehr fein und wird zu größeren Beträgen hin exponentiell gröber. Ob eine Schicht durch die PoT-Rundung stark oder kaum gestört wird, hängt somit weniger von einer festen Bit-Breite ab, wie bei der uniformen Quantisierung, als davon, wie gut die spezifische Verteilung der Gewichtswerte einer Schicht zufällig mit dem PoT-Gitter übereinstimmt. Dieser Punkt ist zentral für das Kernergebnis dieses Projekts (\S\ref{sec:discussion}).

\subsection{Die Hessian-Matrix und Krümmungsanalyse}
\label{sec:fund-hessian}

\textbf{Was die Hessian-Matrix ist.} Für eine Verlustfunktion $L(W)$, die von den Gewichten $W$ eines Netzes abhängt, ist die Hessian-Matrix $H = \nabla^2_W L$ die Matrix der zweiten partiellen Ableitungen des Loss nach jedem Gewichtspaar. Sie beschreibt, wie stark sich die Steigung des Loss ändert, wenn man sich entlang einer Richtung im Gewichtsraum bewegt, anschaulich gesprochen also, wie stark die Verlustlandschaft an der aktuellen Position gekrümmt ist.
\par\medskip

\textbf{Warum sie für Quantisierung relevant ist.} Quantisierung stört die Gewichte um einen kleinen Betrag $\Delta W = W - Q(W)$, wobei $Q(\cdot)$ die Quantisierungsfunktion ist. Eine Taylor-Entwicklung des Loss um den unquantisierten Punkt zweiter Ordnung ergibt
\begin{equation}
\Delta L \approx \nabla L^\top \Delta W + \tfrac{1}{2}\Delta W^\top H \Delta W,
\label{eq:taylor}
\end{equation}
wobei der lineare Term bei einem konvergierten Modell (Gradient nahe null) vernachlässigbar ist, sodass der quadratische, Hessian-gewichtete Term den Loss-Anstieg dominiert. Diese Beziehung ist die theoretische Grundlage dafür, die Krümmung als Proxy für die Quantisierungsempfindlichkeit zu verwenden. Laut Gleichung~\eqref{eq:taylor} sollte eine Schicht mit hoher Krümmung empfindlicher auf eine Störung gleicher Größe reagieren als eine Schicht mit niedriger Krümmung.
\par\medskip

\textbf{Hutchinson-Schätzer für die Spur.} Für moderne Netze mit Millionen Parametern ist die vollständige Hessian-Matrix weder speicherbar noch berechenbar. Praktikabel ist jedoch die Spur $\mathrm{Tr}(H) = \sum_i H_{ii}$, die grob gesprochen die aufsummierte Krümmung über alle Gewichtsrichtungen angibt. Der Hutchinson-Schätzer approximiert die Spur stochastisch. Für zufällige Vektoren $v$ mit $\mathbb{E}[vv^\top] = I$ (üblicherweise Rademacher-Vektoren, deren Komponenten gleichverteilt $\pm 1$ sind) gilt $\mathrm{Tr}(H) = \mathbb{E}[v^\top H v]$. Das Produkt $Hv$, ein Hessian-Vektor-Produkt, lässt sich effizient per automatischer Differentiation berechnen, ohne $H$ jemals explizit zu bilden. Mehrere solcher Zufallsproben werden gemittelt, um die Schätzung zu stabilisieren.
\par\medskip

\textbf{PyHessian.} \citet{yao_pyhessian_2020} implementieren diesen Hutchinson-Ansatz (inklusive Top-Eigenwert-Schätzung über Power-Iteration) als wiederverwendbares Werkzeug für PyTorch-Modelle und etablieren ihn als De-facto-Standard für Krümmungsanalysen in der Quantisierungsliteratur. Dieses Projekt verwendet PyHessian als Grundlage der eigenen Trace-Messungen (\S\ref{sec:methodology}).

\subsection{HAWQ und HAWQ-V2}
\label{sec:fund-hawq}

\citet{dong_hawq_2019} führen HAWQ ein. Dabei werden die Schichten nach ihrem größten Hessian-Eigenwert rangiert, und Schichten mit größerem Eigenwert erhalten mehr Bits in einem gemischt-präzisen Quantisierungsschema. \citet{dong_hawqv2_2020} verfeinern dies zu HAWQ-V2. Statt nur des Eigenwerts wird nun ein Sensitivitäts-Score als Produkt aus Hessian-Spur und quadrierter Störungsgröße gebildet, $\mathrm{Tr}(H_l)\cdot\|\Delta W_l\|^2$ für Schicht $l$, direkt motiviert durch Gleichung~\eqref{eq:taylor}, wobei $\mathrm{Tr}(H_l)$ als Stellvertreter für die Krümmung dient.
\par\medskip

Entscheidend ist, wofür diese Arbeiten ihre Formulierung nutzen und wofür nicht. Beide verwenden den Score, um eine globale Bit-Zuteilung über alle Schichten zu optimieren, und validieren den Erfolg dieser Zuteilung anhand der End-to-End-Genauigkeit des gesamten Netzes nach der Quantisierung. Keine der beiden Arbeiten testet direkt, ob der Score auf Ebene einzelner Schichten tatsächlich mit dem dort auftretenden Schaden korreliert. Ein Netz kann insgesamt eine gute Endgenauigkeit erreichen, obwohl der Score einzelne Schichten falsch einordnet, solange sich die Fehler kompensieren oder die insgesamt zugeteilten Bits ausreichen. Genau diese Lücke, die per-Schicht-Vorhersagegüte des HAWQ-V2-Scores direkt gegen gemessenen Schaden zu testen, adressiert dieses Projekt.

\subsection{Hessian-Spektrum und Random-Matrix-Theorie}
\label{sec:fund-rmt}

Über die Spur hinaus lässt sich das Spektrum, also die Menge der Eigenwerte, der Hessian-Matrix untersuchen. \citet{sagun_singularity_2016, sagun_empirical_2017} zeigen empirisch, dass das Hessian-Spektrum tiefer neuronaler Netze typischerweise aus zwei Teilen besteht, einem Bulk nahe null, der die meisten Eigenwerte enthält, und wenigen großen Ausreißer-Eigenwerten, die mit der Anzahl der Klassen des Klassifikationsproblems zusammenhängen. Diese Bulk-Ausreißer-Trennung motiviert die Vorstellung, dass ein Netz effektiv nur entlang weniger Richtungen scharf gekrümmt ist.
\par\medskip

\citet{karakida_fisher_2019} liefern hierzu eine theoretische Grundlage aus der Random-Matrix-Theorie (RMT). Für zufällig initialisierte, untrainierte Netze lässt sich die Fisher-Information, die eng mit der Hessian-Matrix nahe einem Minimum verwandt ist, im Mittelfeld-Limit vorhersagen, inklusive einer charakteristischen Tiefenabhängigkeit. Diese Vorhersage ist der theoretische Bezugspunkt für die Random-Init-Kontrolle dieses Projekts (Phase~2, \S\ref{sec:phase2}). Sie erlaubt zu unterscheiden, ob eine beobachtete Krümmungsspitze in einer bestimmten Schicht bereits bei Initialisierung vorhanden ist (\enquote{architektonisch}, durch die Netzstruktur selbst erklärbar) oder erst durch Training entsteht (\enquote{gelernt}).
\par\medskip

Eine weitere Zerlegungsidee ist die Kronecker-faktorisierte Näherung der Krümmung (KFAC, \citet{martens_kfac_2015}) sowie ihre Erweiterung auf Faltungsschichten (KFC, \citet{grosse_kfc_2016}). Die Grundidee besteht darin, die Fisher-Information einer einzelnen Schicht näherungsweise als Kronecker-Produkt zweier kleinerer Matrizen zu faktorisieren, $F \approx A \otimes G$, wobei $A$ die Kovarianz der Schicht-Eingaben und $G$ die Kovarianz der Verlust-Gradienten bezüglich der Schicht-Ausgaben erfasst. Gelingt diese Faktorisierung, so ließe sich die Spur einer Schicht in einen Eingabe-getriebenen und einen Gradienten-getriebenen Anteil zerlegen, was eine potenziell aufschlussreiche Diagnose dafür wäre, warum eine Schicht hohe Krümmung aufweist. Sowohl die RMT-Vorhersage als auch die KFAC-Zerlegung waren für dieses Projekt vielversprechende Werkzeuge, um die architektonischen Ursachen beobachteter Krümmungsspitzen zu erklären (Phase~6, \S\ref{sec:phase6}). Wie in \S\ref{sec:results} gezeigt wird, führte insbesondere die KFAC-Zerlegung im PoT-Kontext jedoch nicht zum Ziel, da die vorhergesagte und die gemessene Spur um mehrere Größenordnungen voneinander abweichen.

\subsection{Aktivierungsquantisierung und Range-Ausreißer}
\label{sec:fund-activation}

Alle bisherigen Konzepte betreffen die Quantisierung von Gewichten. Ein separates, in der Literatur gut dokumentiertes Problem ist die Quantisierung von Aktivierungen, also der Zwischenausgaben einer Schicht während der Inferenz. Aktivierungen folgen einem anderen Schadensmechanismus als Gewichte. Statt gleichmäßig verteilter Fehler treten oft wenige, aber extreme Ausreißer-Kanäle auf, deren Wertebereich um Größenordnungen über dem Rest liegt. Ein einzelner Skalierungsfaktor pro Tensor muss diesen gesamten Bereich abdecken, was die effektive Auflösung für alle übrigen, unauffälligen Kanäle drastisch reduziert. \citet{xiao_smoothquant_2023} (SmoothQuant) und \citet{lin_awq_2024} (AWQ) zeigen dies für große Sprachmodelle und schlagen vor, die Schwierigkeit vor der Quantisierung von den Aktivierungen auf die leichter quantisierbaren Gewichte zu verschieben. \citet{li_brecq_2021} (BRECQ) verfolgt einen alternativen, blockweisen Rekonstruktionsansatz für PTQ, der Aktivierungs- und Gewichtsquantisierung gemeinsam über kleine Netzblöcke optimiert.
\par\medskip

Für dieses Projekt ist relevant, dass Aktivierungs-Ausreißer kanalweise auftreten und primär von der Eingabeverteilung bestimmt werden, während die Hessian-Krümmung eine Eigenschaft der Gewichte relativ zum Loss ist. Es gibt keinen a-priori-Grund, warum ein Hessian-basierter Score, der für Gewichtsschaden entwickelt wurde, auch Aktivierungsschaden vorhersagen sollte. Diesen Punkt greift die Dekomposition in Gewichts- und Aktivierungsanteile (Phase~1, \S\ref{sec:phase1}) empirisch auf.